\documentclass[12pt]{article}
\usepackage[utf8]{inputenc}
\usepackage[T1]{fontenc}
\usepackage{hyperref}
\title{Yaroam\"e Language Resources for K-12 Curriculum Development}
\author{Compiled by Claude Sonnet 5.5 (claude-sonnet-5-5) for Andr\'as Kornai}
\date{}
\begin{document}
\maketitle

\section{Introduction}
Yaroam\"e (ISO 639-3 yro) is a Yanomaman language of Brazil. An Iphan/Instituto Socioambiental article reports 371 speakers and describes it as one of the Yanomami languages in a concerning situation. I found no dictionary or school materials specific to Yaroam\"e in the searches. The materials below concern the Yanomami languages in general or other Yanomaman languages, so they should not be treated as Yaroam\"e resources without checking.

\section{Primary Educational Resources}
No Yaroam\"e primers or textbooks were found. The following are about Yanomami schooling in general:
\begin{itemize}
\item Thesis \emph{Educa\c{c}\~ao escolar yanomami e potiguara}, UFPB: \url{https://repositorio.ufpb.br/jspui/handle/tede/6455}
\item Brazilian federal prosecutors' (MPF) call to restructure schooling for Yanomami children and youth in Roraima: \url{https://www.mpf.mp.br/o-mpf/unidades/pr-rr/noticias/mpf-pede-reestruturacao-urgente-da-educacao-escolar-para-3-5-mil-criancas-e-jovens-yanomami-em-roraima}
\end{itemize}

\section{Linguistic and Reference Materials}
\begin{itemize}
\item Iphan recognises Yanomami linguistic diversity (Instituto Socioambiental): \url{https://site-antigo.socioambiental.org/en/print/6272}
\item Fargetti, illustrated Yanomami--Spanish dictionary (Mana, 2008), UNESP record. It is not specific to Yaroam\"e: \url{https://acervodigital.unesp.br/handle/11449/6410}
\item NSF grant for a Ninam (another Yanomaman language) dictionary: \url{https://grantome.com/grant/NSF/BCS-0963595}
\end{itemize}

\section{Religious and Community Resources}
None were found.

\section{Digital and Audio Resources}
None were found.

\section{Curriculum Development Context}
The language is thought to be spoken by the whole community. Literacy materials and teacher training in Yaroam\"e itself appear to be lacking. The sources cite a shortage of written Yanomami material as an obstacle to schooling. Any curriculum would need to be built with Yanomami communities and organisations.
\end{document}